\documentclass[12pt]{article}

\usepackage[a4paper, margin=2.5cm]{geometry}
\usepackage{amsmath}
\usepackage{amssymb}
\usepackage{amsthm}
\usepackage[colorlinks=true, linkcolor=blue, citecolor=blue, urlcolor=blue]{hyperref}
\usepackage{booktabs}
\usepackage{natbib}

\newtheorem{theorem}{Theorem}
\newtheorem{proposition}[theorem]{Proposition}
\newtheorem{corollary}[theorem]{Corollary}
\newtheorem{definition}[theorem]{Definition}
\newtheorem{remark}[theorem]{Remark}
\newtheorem{lemma}[theorem]{Lemma}
\newtheorem{conjecture}[theorem]{Conjecture}
\newtheorem{observation}[theorem]{Observation}

\newcommand{\phig}{\varphi}
\newcommand{\Stab}{\mathrm{Stab}}

\title{Cross-$\beta$ Amyloid Fibril Twist from $2I$ Helical-Orbit Classes\\[6pt]
\large A Commensurability Prediction for Diabetes, Alzheimer's, Parkinson's, Prion, and Cancer-Associated Aggregation}
\author{%
  Lee Smart\\[4pt]
  \textit{Institute of Vibrational Field Dynamics}\\[2pt]
  \texttt{contact@vibrationalfielddynamics.org}\\[2pt]
  \texttt{@vfd\_org}%
}
\date{April 2026}

\begin{document}

\maketitle

\noindent\textbf{Status:} Preprint (not peer-reviewed)

\begin{abstract}
Cross-$\beta$ amyloid fibrils are the shared pathological
architecture of a wide disease family: type 2 diabetes (islet
amyloid polypeptide, IAPP), Alzheimer's disease (amyloid-$\beta$,
tau), Parkinson's disease ($\alpha$-synuclein), prion diseases
(PrP$^{\mathrm{Sc}}$), systemic amyloidoses, and cancer-associated
gain-of-function aggregation (oncogenic p53). Each fibril species
carries a small per-strand twist angle $\theta_{\mathrm{twist}}$
(typically $1^{\circ}$--$5^{\circ}$) accumulating along the fibril
axis; per-strand stacking is fixed near $d_\beta = 4.7$ \AA\ by
$\beta$-sheet hydrogen-bond geometry.

We formalise (from the Vibrational Field Dynamics cascade
Life-rung machinery) a discrete set of \emph{cascade-admissible}
per-strand twists, obtained by:

\begin{itemize}
\item \textbf{Lemma~\ref{lem:L1}} identifying the preimage in
  $2I \subset \mathrm{SU}(2)$ of a 5-fold stabiliser in
  $I \cong A_5$ as the cyclic group $C_{10}$ of order 10 by the
  Schur--Zassenhaus theorem (the extension splits because the
  kernel and quotient orders $2$ and $5$ are coprime). The
  smallest nonzero single-$2I$-element rotation about a 5-fold
  axis is $36^{\circ}$;
\item \textbf{Definition~\ref{def:D2}} parametrising
  commensurable-twist families $\theta(p, q, N) = 36^{\circ}
  (p - 10q)/N$ with $N$ an explicit integer resolution
  parameter; and
\item \textbf{Definition~\ref{def:D3}} imposing an exact
  amyloid-scale filter: $M \cdot \theta = 360^{\circ}$ with
  integer $M$ and crossover length $M \cdot d_\beta \in [100, 500]$
  \AA, giving $M \in [22, 106]$ and
  $\theta \in [3.3962^{\circ}, 16.3636^{\circ}]$.
\end{itemize}

An enumeration script produces the $24$-element admissible set
for resolution $N \le 20$; all rows satisfy $M \theta = 360^{\circ}$
exactly. The resulting set contains the $\theta = 36^{\circ}/N$
ladder for $N = 3, \ldots, 10$ (i.e.
$\{3.6^{\circ}, 4.0^{\circ}, 4.5^{\circ}, 5.14^{\circ}, 6.0^{\circ},
7.2^{\circ}, 9.0^{\circ}, 12.0^{\circ}\}$) plus additional
$(p, q, N)$-combinations such as $\{3.79^{\circ}, 4.24^{\circ},
\ldots\}$.

\textbf{Empirical conjecture and falsifier.} The central prediction
(Conjecture~\ref{conj:C1}) is that observed cross-$\beta$ fibril
per-strand twists cluster on this 24-element admissible set
rather than distributing uniformly over $[3.4^{\circ}, 16.4^{\circ}]$.
Ultra-slow-twist polymorphs ($\theta < 3.4^{\circ}$, e.g.\ tau PHF
$\sim 1^{\circ}$) lie outside the current filter and are flagged
as an open regime (Remark~\ref{rem:R2}). The empirical test is
pending curation of a cryo-EM amyloid-twist dataset.

No empirical comparison is carried out in this paper; the
comparison pipeline follows the WO-1 capsid
template~\cite{SmartCapsid} with \texttt{wo2\_amyloid\_fetch.py}
and \texttt{wo2\_compare.py} planned in the same
\texttt{WAITING\_FOR\_DATA} style. Every claim is classified as
\textbf{derived}, \textbf{identified}, or \textbf{open}.
\end{abstract}

%==================================================================
\section{Introduction and Scope}\label{sec:intro}
%==================================================================

\subsection{Cross-$\beta$ architecture}

A cross-$\beta$ amyloid fibril is a linear polymer of
$\beta$-strands stacked perpendicular to a common fibril axis,
with the strands held together by main-chain hydrogen bonds. Key
geometric parameters:

\begin{itemize}
\item \textbf{Strand stacking}: $d_\beta \approx 4.7$ \AA\ per
  strand (fixed by canonical $\beta$-sheet hydrogen-bonding).
\item \textbf{Per-strand twist angle} $\theta_{\mathrm{twist}}$:
  a small rotation about the fibril axis, typically
  $1^{\circ}$--$5^{\circ}$ per strand.
\item \textbf{Crossover distance}:
  $L_{\mathrm{co}} = 360^{\circ} \cdot d_\beta / \theta_{\mathrm{twist}}$;
  observed range roughly $100$--$500$ \AA, with some polymorphs
  extending to $> 1000$ \AA.
\item \textbf{Protofilament count}: typically $1$--$4$ with
  $C_n$ or $D_n$ bundle symmetry.
\end{itemize}

Cross-$\beta$ architecture is shared across the amyloid-disease
family. The present paper asks: does the cascade Life-rung
substrate constrain $\theta_{\mathrm{twist}}$ to a discrete
admissible set, or is $\theta_{\mathrm{twist}}$ set purely by
polypeptide backbone mechanics with no substrate-level
constraint?

\subsection{Disease relevance}

Amyloid fibril formation is central to:

\begin{itemize}
\item \textbf{Type 2 diabetes}: pancreatic $\beta$-cell
  dysfunction is associated with islet amyloid polypeptide
  (IAPP / amylin) fibril deposition.
\item \textbf{Alzheimer's disease}: amyloid-$\beta$ (A$\beta$)
  plaques and tau neurofibrillary tangles, both cross-$\beta$.
\item \textbf{Parkinson's disease}: Lewy bodies are
  $\alpha$-synuclein fibrils with cross-$\beta$ architecture.
\item \textbf{Prion diseases}: PrP$^{\mathrm{Sc}}$ cross-$\beta$
  aggregates.
\item \textbf{Cancer-associated aggregation}: some oncogenic
  missense mutants of p53 form amyloid-like fibrils as a
  gain-of-function mechanism.
\end{itemize}

A prediction constraining fibril twist across all of these is
simultaneously a prediction for every disease in the list.

\subsection{Relation to sibling WOs}

This paper is WO-2 of the biology-rung programme
\cite{SmartCapsid,SmartPhyllotaxis,SmartTubulin,SmartBioelectric}.
It depends on:
\begin{itemize}
\item WO-1 (capsid) for the $2I / A_5$ substrate and the
  catalogue convention.
\item The upstream B2 decagram result
  \cite[\S5.1]{CascadeBio}: the natural fast helical orbit on a
  5-fold axis of the 600-cell is the decagram, with $10$ vertices
  per $2\pi$-turn ($36^{\circ}$/step), matching DNA's 10 bp/turn.
\end{itemize}
The present paper extends B2 from the single fast decagram pitch
to the full commensurable-twist family.

%==================================================================
\section{Helical-Orbit Structure on the 5-Fold Axis}\label{sec:orbits}
%==================================================================

\subsection{The $C_{10}$ stabiliser}

\begin{lemma}[$\blacktriangleright$ Admissible angles on a 5-fold axis]\label{lem:L1}
Let $A$ be a 5-fold axis of the 600-cell. The preimage of
$\Stab_I(A) \cong C_5$ in the double cover $2I \to I$ is cyclic
of order $10$. The SU(2) rotation angles (in the convention
$R = \exp(i(\alpha/2)\hat{n} \cdot \hat{\sigma})$ for axis unit
vector $\hat{n}$) of the $10$ elements are
\[
\Stab_{2I}(A) = \{\alpha : \alpha = 2\pi k / 10, k = 0, 1, \ldots, 9\}.
\]
In degrees: $\{0^{\circ}, 36^{\circ}, 72^{\circ}, 108^{\circ},
144^{\circ}, 180^{\circ}, 216^{\circ}, 252^{\circ}, 288^{\circ},
324^{\circ}\}$. The smallest nonzero single-$2I$-element
rotation about $A$ is $36^{\circ}$.
\end{lemma}

\begin{proof}
The short exact sequence $1 \to \{\pm 1\} \to 2I \to I \to 1$
has kernel of order $2$ and projects $\Stab_I(A) \cong C_5$ onto
the 5-fold stabiliser. The preimage is a central extension of
$C_5$ by $C_2$. Since $\gcd(2, 5) = 1$, the
Schur--Zassenhaus theorem gives a splitting; explicitly the
preimage is $C_5 \times C_2 \cong C_{10}$ (using the isomorphism
of finite abelian groups with coprime factorisation). The
rotation angles of $C_{10}$'s ten elements as SU(2) rotations
about $A$ are the $10$th roots of unity, i.e.\ $2\pi k / 10$ for
$k = 0, \ldots, 9$.
\end{proof}

\begin{remark}[$\triangleright$ The small-angle problem]\label{rem:R1}
Empirical amyloid twists cluster in the $1^{\circ}$--$5^{\circ}$
range. By Lemma~\ref{lem:L1} these cannot be achieved as single-
$2I$-element rotations about a 5-fold axis (minimum is
$36^{\circ}$). The small-angle twist must therefore emerge from
\emph{compounded} group elements, not a single $R_\alpha$: either
commensurable residues of multi-step helical walks, or orbits
restricted to subgroups of $2I$, or cell-level (not vertex-level)
orbits. This document focuses on the compound-residue route.
\end{remark}

\subsection{Commensurable twist family}

\begin{definition}[$\blacktriangleright$ Commensurable twist]\label{def:D2}
Fix $\alpha \in \Stab_{2I}(A) \setminus \{0\}$, an integer
$N \ge 1$, and integers $(p, q)$ with $\gcd(p, q, N) = 1$. The
\emph{commensurable twist at $(\alpha; p, q, N)$} is
\[
\theta(\alpha; p, q, N) := (p \alpha - q \cdot 2\pi) / N
\qquad (\text{reduced mod } 2\pi).
\]
$N$ is an explicit parameter of the definition (not derived from
$(p, q, \alpha)$). Geometrically: after $N$ strands along the
fibril, the helical pattern has accumulated $N\theta$ of
in-plane rotation, which equals $p$ applications of $R_\alpha$
minus $q$ full turns.

For $\alpha = 36^{\circ}$, one computes
$\theta = 36^{\circ} \cdot (p - 10q)/N$; setting $s := p - 10q$,
every commensurable twist is a rational multiple of
$36^{\circ}/N$ within the $36^{\circ}/N \cdot \mathbb{Z}$ ladder.
Other $\alpha \in \Stab_{2I}(A)$ (multiples of $36^{\circ}$) give
no angles outside this ladder.
\end{definition}

\begin{lemma}[$\blacktriangleright$ Countability and zero-accumulation]\label{lem:L2}
The union $\cup_{N \ge 1} \Theta_{\mathrm{adm}}(A; N)$ is a
countable subset of $\{36^{\circ} \cdot m/N : m, N \in \mathbb{Z},
N \ge 1\}$ and has $0$ as an accumulation point.
\end{lemma}

\begin{proof}
Every commensurable twist is of the form $36^{\circ}(p-10q)/N =
36^{\circ}m/N$ with $m = p - 10q \in \mathbb{Z}$; this is the
rational-number ladder, which is countable. For accumulation at
zero: given $\epsilon > 0$, choose any $N > 36^{\circ}/\epsilon$
and $m = 1$; then $\theta = 36^{\circ}/N < \epsilon$.
\end{proof}

\subsection{Amyloid-scale filter}

\begin{definition}[$\blacktriangleright$ Amyloid-scale compatibility]\label{def:D3}
A twist $\theta$ is \emph{amyloid-scale compatible under
$n$-protofilament bundle symmetry} if there exists a positive
integer $M$ with
\[
M \cdot \theta \in \frac{360^{\circ}}{n}\,\mathbb{Z}, \qquad
\text{i.e. } M\theta\text{ is an exact multiple of } 360^{\circ}/n
\]
(the strict full-turn case used in this definition is $n = 1$;
the generalised $n \ge 2$ case accommodates multi-protofilament
$C_n$-bundles whose cross-section repeats every $360^{\circ}/n$
and is enumerated in build $B_{\mathrm{crossover}}$ at
\texttt{scripts/wo2\_sim\_close\_crossover.py}), \textbf{and} the
crossover length satisfies
\[
M \cdot d_\beta = M \cdot 4.7 \text{ \AA} \in [100 \text{ \AA}, 500 \text{ \AA}].
\]
Solving: $M \in \{22, 23, \ldots, 106\}$ and therefore
\[
\theta \in [360^{\circ}/106, 360^{\circ}/22] = [3.3962^{\circ}, 16.3636^{\circ}].
\]
The $M$-integrality requirement is the substantive constraint
beyond the bare crossover length range: it picks out
commensurable-twist family members whose fibril crossovers close
\emph{exactly} modulo $360^{\circ}$, excluding twists that only
close approximately.
\end{definition}

\begin{remark}[$\triangleright$ Ultra-slow-twist polymorphs]\label{rem:R2}
Some tau PHF and $\alpha$-synuclein fibril polymorphs report
per-strand twists below $3.4^{\circ}$, corresponding to
crossover distances $> 1500$ \AA. These lie outside
Definition~\ref{def:D3}'s filter. Three candidate explanations
are recorded (not resolved here): (i) they sit at higher
resolution $N > 106$, technically within the commensurable set
but outside the canonical short-crossover filter; (ii) they
correspond to cell-level orbits (rather than vertex-level) on
the 600-cell; (iii) they correspond to different axis types
(3-fold or 2-fold rather than 5-fold). An open sub-problem is
to compute the cell-level orbit spectrum and test whether it
absorbs these polymorphs.
\end{remark}

%==================================================================
\section{Enumeration and Predictions}\label{sec:enum}
%==================================================================

\subsection{The 24-element admissible set}

The admissibility condition $M = 10N/s$ (integer, in
$[22, 106]$) with $\theta = 36^{\circ} s/N$ and $(s, N)$ coprime
gives a finite set for each resolution bound $N_{\max}$. For
$N_{\max} = 20$, the script
\texttt{papers/biology-rung/scripts/wo2\_helical\_orbits.py}
produces $24$ distinct admissible $\theta$-values, all
satisfying $M\theta = 360^{\circ}$ exactly. A representative
subset is shown in Table~\ref{tab:admissible}.

\begin{table}[h]
\centering
\caption{Representative cascade-admissible amyloid per-strand
twists (resolution $N \le 20$), with exact integer crossover $M$
and resulting crossover length. Full 24-element list in
\texttt{data/wo2\_prediction.csv}.}
\label{tab:admissible}
\begin{tabular}{@{}rrrrr@{}}
\toprule
$\theta$ (deg) & $s$ & $N$ & $M$ & $L_{\mathrm{co}}$ (\AA) \\
\midrule
3.6000 & 1 & 10 & 100 & 470.0 \\
4.0000 & 1 & 9 & 90 & 423.0 \\
4.5000 & 1 & 8 & 80 & 376.0 \\
5.1429 & 1 & 7 & 70 & 329.0 \\
6.0000 & 1 & 6 & 60 & 282.0 \\
7.2000 & 1 & 5 & 50 & 235.0 \\
9.0000 & 1 & 4 & 40 & 188.0 \\
12.0000 & 1 & 3 & 30 & 141.0 \\
14.4000 & 2 & 5 & 25 & 117.5 \\
16.3636 & 5 & 11 & 22 & 103.4 \\
\bottomrule
\end{tabular}
\end{table}

\subsection{Empirical conjecture}

\begin{conjecture}[$\circ$ Amyloid-twist selection rule]\label{conj:C1}
Observed cross-$\beta$ amyloid per-strand twists cluster on the
enumerated admissible set
$\Theta_{\mathrm{adm}}(A) \cap [3.3962^{\circ}, 16.3636^{\circ}]$
rather than distributing uniformly over that range.

\textbf{Falsifier.} A uniform distribution of observed
$\theta_{\mathrm{twist}}$ over $[3.4^{\circ}, 16.4^{\circ}]$
across a curated cryo-EM amyloid dataset, with no clustering at
the admissible-set values, would falsify this conjecture.
\textbf{Status}: untested; comparison script planned following
the WO-1 WAITING\_FOR\_DATA template.
\end{conjecture}

\begin{conjecture}[$\circ$ Chirality bias; qualitative]\label{conj:C2}
The $2I$ group carries intrinsic chirality per
\cite[\S 2.7, \S B5]{CascadeBio}. Cross-$\beta$ fibrils may show
a handedness bias inherited from the cascade substrate; upstream
work only establishes that $2I$-orbits are intrinsically chiral
and that a god-prime $2^n + 1$ argument would select one
enantiomer, but the \emph{direction} of the bias is not
derived. The claim here is qualitative: a nonzero asymmetry
may be present. Magnitude and sign are open.
\end{conjecture}

\begin{conjecture}[$\circ$ Protofilament count from $2I$ subgroups]\label{conj:C3}
Admissible protofilament bundle symmetries correspond to
non-trivial cyclic and binary-dihedral subgroups of $2I$
(principally $C_1, C_2, C_3, C_5, D_2$, etc.), giving the
1--4-protofilament bundle range commonly observed in cryo-EM.
The absence of, e.g., $C_7$-symmetric protofilament bundles
across a careful cryo-EM survey would be consistent with (but
not necessarily required by) $2I$ subgroup exhaustion.
\end{conjecture}

\subsection{Per-disease empirical targets}

The predictions of Conjecture~\ref{conj:C1} apply uniformly
across the amyloid-disease family. Disease-specific comparison
targets (all pending data curation):
\begin{itemize}
\item IAPP (type 2 diabetes): polymorph-dependent, some cryo-EM
  structures reporting $\theta \approx 4.5^{\circ}$ which would
  match the admissible $\theta = 36^{\circ}/8$ entry.
\item A$\beta_{1-42}$ (Alzheimer's): polymorphs reported in
  $1^{\circ}$--$5^{\circ}$ range; many in Remark~\ref{rem:R2}'s
  open regime.
\item Tau PHF (Alzheimer's): $\sim 1^{\circ}$, in
  Remark~\ref{rem:R2}'s open regime.
\item $\alpha$-synuclein (Parkinson's): $\sim 2^{\circ}$, in
  Remark~\ref{rem:R2}'s open regime.
\item PrP$^{\mathrm{Sc}}$ (prion): structure-dependent.
\item p53-cancer aggregation: limited structural data; flagged
  as an interesting stretch target for future cryo-EM work.
\end{itemize}
The sub-3.4$^{\circ}$ polymorphs fall outside the present
filter; whether they are absorbed by cell-level orbits
(Remark~\ref{rem:R2}) is an explicit open question, not resolved
in this paper.

%==================================================================
\section{Summary and Logical Chain}\label{sec:chain}
%==================================================================

\paragraph{What is derived ($\blacktriangleright$).}
\begin{itemize}
\item Lemma~\ref{lem:L1}: $\Stab_{2I}(A) \cong C_{10}$ on a
  5-fold axis, smallest nonzero rotation $36^{\circ}$.
\item Definition~\ref{def:D2}: commensurable-twist family with
  explicit resolution parameter $N$ (non-circular).
\item Lemma~\ref{lem:L2}: countable set with zero-accumulation.
\item Definition~\ref{def:D3}: exact-$M$ amyloid-scale filter,
  $\theta \in [3.3962^{\circ}, 16.3636^{\circ}]$.
\item Computational enumeration: 24 admissible $\theta$-values
  at resolution $N \le 20$, all satisfying $M\theta = 360^{\circ}$
  exactly (script at
  \texttt{papers/biology-rung/scripts/wo2\_helical\_orbits.py}).
\end{itemize}

\paragraph{What is identified ($\triangleright$).}
\begin{itemize}
\item Remark~\ref{rem:R1}: small-angle amyloid twists cannot be
  single-$2I$-element rotations; must emerge from compound steps.
\item Remark~\ref{rem:R2}: ultra-slow-twist polymorphs
  ($< 3.4^{\circ}$) lie outside the present filter; three
  candidate resolutions flagged as open.
\item The decagram substrate of B2 \cite[\S 5.1]{CascadeBio} is
  the 5-fold-axis origin; other axis types (3-fold, 2-fold) not
  enumerated here.
\end{itemize}

\paragraph{What is open ($\circ$).}
\begin{itemize}
\item Conjecture~\ref{conj:C1} (amyloid-twist selection):
  empirical test pending curated cryo-EM dataset.
\item Conjecture~\ref{conj:C2} (chirality bias): qualitative
  only; magnitude and sign not derived.
\item Conjecture~\ref{conj:C3} (protofilament count from
  subgroups): conjecture, not proof.
\item Cell-level orbit enumeration for sub-3.4$^{\circ}$
  polymorphs (Remark~\ref{rem:R2}).
\item Multi-axis extension (3-fold and 2-fold axes of the
  600-cell).
\item Whether the amyloid selection rule is the spectrum of a
  closure operator parallel to WO-1's $\mathcal{T}_{\mathrm{casc}}$
  \cite[T.1]{SmartCapsid}.
\end{itemize}

%==================================================================
\section{Reproducibility}\label{sec:repro}
%==================================================================

This manuscript lives at
\texttt{papers/biology-rung-amyloid/biology-rung-amyloid.tex};
the mathematical source is \texttt{derivation-amyloid.md} in
\texttt{papers/biology-rung/}, the enumeration script is
\texttt{scripts/wo2\_helical\_orbits.py}, and the admissible
CSV output is \texttt{data/wo2\_prediction.csv}.

The empirical-comparison scripts
(\texttt{wo2\_amyloid\_fetch.py},
\texttt{wo2\_compare.py}) are planned following the WO-1
WAITING\_FOR\_DATA template but not yet written. A cryo-EM
amyloid-twist dataset is pending curation.

\begin{thebibliography}{99}

\bibitem{CascadeBio}
\emph{Cascade Biology Rung: The 600-Cell Cell-Level Substrate},
internal manuscript,
\texttt{papers/cascade-derivation/cascade-bio.md}. Preprint,
Institute of Vibrational Field Dynamics.

\bibitem{SmartCapsid}
Smart, L.,
\emph{Caspar--Klug Triangulation Numbers from a Face-Level
Cascade Operator} (WO-1),
\texttt{papers/biology-rung-capsid/biology-rung-capsid.tex}.
Preprint, Institute of Vibrational Field Dynamics, 2026.

\bibitem{SmartPhyllotaxis}
Smart, L.,
\emph{Fibonacci Phyllotaxis from the 600-Cell's $\phig$-Helical
Closure Orbits} (WO-5),
\texttt{papers/biology-rung-phyllotaxis/biology-rung-phyllotaxis.tex}.
Preprint, Institute of Vibrational Field Dynamics, 2026.

\bibitem{SmartTubulin}
Smart, L.,
\emph{The 13-Protofilament Microtubule as a $C_{13}$ Selector}
(WO-3),
\texttt{papers/biology-rung-tubulin/biology-rung-tubulin.tex}.
Preprint, Institute of Vibrational Field Dynamics, 2026.

\bibitem{SmartBioelectric}
Smart, L.,
\emph{Bioelectric Prepattern Eigenmodes as Morphogenetic
Attractors} (WO-4),
\texttt{papers/biology-rung-bioelectric/biology-rung-bioelectric.tex}.
Preprint, Institute of Vibrational Field Dynamics, 2026.

\end{thebibliography}

\end{document}
